% !TeX root = ../main.tex

\xchapter{相关理论基础与技术原理}{Relevant theoretical foundations and technical principles}

本章围绕本文所研究的免训练视频摘要问题，系统介绍后续方法设计所依赖的理论基础与关键技术原理。首先，结合本文研究对象给出视频摘要任务的形式化定义，并说明关键帧、关键片段与训练自由设定等基础概念；随后，针对第三章所采用的基于大语言模型的层级语义视频摘要方法，介绍 Transformer、自提示上下文学习以及跨模态语义对齐等核心原理；接着，针对第四章所采用的多智能体协同推理框架，阐述基于大语言模型的 Agent 推理机制与多 Agent 协作机制；最后，对全文统一使用的视频摘要评测指标与评测协议进行说明，为后续章节的方法描述与实验分析提供理论支撑。

\xsection{视频摘要任务建模与基本概念}{Task modeling and basic concepts of video summarization}

视频摘要旨在从原始视频中自动选择少量最具代表性、信息量和主题相关性的内容片段，以在显著压缩时长的同时尽可能保留原视频的核心语义与叙事结构。与静态图像摘要不同，视频摘要不仅需要判断单帧画面的视觉显著性，还需要综合考虑事件连续性、跨时间依赖、主题一致性以及冗余抑制等因素。因此，该任务本质上是一个同时涉及时序建模、语义理解与组合优化的结构化决策问题\cite{zhang2016video,rochan2018video,apostolidis2022summarizing}。

从问题属性上看，视频摘要至少包含四个彼此耦合的核心难点。第一，视频内容具有显著的时间冗余，相邻帧往往在视觉上高度相似，但在语义层面又可能承载不同的事件阶段，因此系统既要抑制重复信息，又不能破坏事件完整性。第二，视频摘要具有较强的主题依赖性，同一片段在不同视频主题下的重要程度可能完全不同，例如在教程类视频中某个操作步骤可能极为关键，而在风景展示类视频中相似画面则未必应被保留。第三，视频摘要具有显著的主观性，不同观看者对“重要内容”的判断往往受叙事偏好、信息需求和视觉注意力差异影响。第四，最终摘要还必须满足长度预算约束，这使其不仅是一个排序问题，也是一个受约束的子集选择问题。

从输出形式看，视频摘要通常可以分为关键帧摘要与关键片段摘要两类。关键帧摘要以若干离散帧构成压缩结果，适合快速浏览视频内容；关键片段摘要则以时间连续的镜头或片段作为最小选择单元，更符合真实观看体验，也是 SumMe 和 TVSum 等主流基准中更常见的评测形式。无论采用何种输出形式，其核心目标都可以归结为：在给定长度约束下，从原视频中选出一组最能代表整体语义且冗余最小的内容单元。

设输入视频记为
\begin{equation}
\mathcal{V}=\{f_t\mid t=1,2,\ldots,T\},
\end{equation}
式中，$\mathcal{V}$——输入视频序列；$f_t$——第 $t$ 帧；$T$——视频总帧数。视频摘要系统的目标是在长度预算约束下，生成一个与原视频对应的帧级重要性分数序列 $\mathbf{s}=(s_1,s_2,\ldots,s_T)$，或进一步生成二值选择向量
\begin{equation}
\mathbf{y}=\{y_t\mid y_t\in\{0,1\},\ t=1,2,\ldots,T\},
\end{equation}
式中，$s_t$——第 $t$ 帧的重要性分数；$y_t$——第 $t$ 帧是否被选入摘要的二值决策变量。当 $y_t=1$ 时表示该帧对应内容被保留，反之则被舍弃。若将连续帧进一步组织为片段集合 $\mathcal{S}=\{S_1,S_2,\ldots,S_N\}$，则视频摘要问题也可表示为在预算约束下选择最优片段子集：
\begin{equation}
\max_{\mathbf{z}} \sum_{k=1}^{N} z_k u_k, \qquad
\text{s.t.}\ \sum_{k=1}^{N} z_k l_k \leq B,\ z_k\in\{0,1\},
\end{equation}
式中，$z_k$——第 $k$ 个片段的选择变量；$u_k$——第 $k$ 个片段的重要性收益；$l_k$——片段长度；$B$——摘要总长度预算。该形式清晰表明，视频摘要同时包含重要性估计与长度受限选择两个核心环节。

围绕上述形式化定义，本文主要使用以下基础概念。其一，关键帧是指能够代表局部场景或事件语义的静态帧，适合表达视频中的代表性瞬间；关键片段则由一段时间连续的帧组成，能够更完整地保留动作发展、事件过渡与上下文关系。其二，帧级重要性分数是系统对每一时刻内容价值的连续估计，通常用于后续平滑、排序与摘要生成；片段级重要性分数则是对时间区间整体价值的评价，更适合与场景分割或固定窗口分段策略结合。其三，冗余抑制是视频摘要任务区别于一般排序任务的重要特征，即高分内容若在语义上高度重复，仍不应被大量保留，否则会削弱摘要的信息密度与观看效率。

为了更准确地刻画视频摘要的目标，许多方法会在重要性收益之外显式考虑代表性与多样性。若记片段间的语义相似度为 $\mathrm{sim}(S_i,S_j)$，则摘要目标可进一步扩展为
\begin{equation}
\max_{\mathbf{z}} \sum_{k=1}^{N} z_k u_k - \lambda \sum_{i=1}^{N}\sum_{j=1}^{N} z_i z_j\, \mathrm{sim}(S_i,S_j),
\end{equation}
式中，$\lambda$——冗余惩罚系数；$\mathrm{sim}(S_i,S_j)$——片段 $S_i$ 与 $S_j$ 的语义相似度。该形式说明，理想摘要不仅应包含高价值内容，还应避免选择大量彼此相似的片段。对本文而言，这一思想在第三章体现为全局主题约束下的语义一致性建模，在第四章则体现为多智能体对重复片段和过渡片段的差异化判断。

在公开研究中，视频摘要通常基于 SumMe 和 TVSum 两个基准数据集进行评测。SumMe 更强调开放场景下的主观摘要偏好，视频内容类型多样、标注差异较大；TVSum 则覆盖新闻、教程、运动、纪录片等多种网络视频类型，并提供逐帧重要性评分，更适合分析模型在连续排序层面的性能\cite{gygli2014creating,song2015tvsum}。这两个数据集共同构成了本文实验部分的主要评价基础，也使第二章中对任务定义和评测协议的讨论具有统一参照系。

本文关注的训练自由视频摘要可进一步形式化为：在给定基础模型参数 $\theta$ 固定的条件下，通过构造任务提示 $\pi$、上下文表示 $\mathcal{C}$ 与语义映射函数 $g_{\theta}(\cdot)$，直接得到片段或帧的重要性估计
\begin{equation}
s_t = g_{\theta}(f_t,\mathcal{C}_t,\pi), \qquad \theta\ \text{fixed},
\end{equation}
式中，$g_{\theta}(\cdot)$——固定参数模型在给定上下文和提示下的推理映射；$\mathcal{C}_t$——与第 $t$ 帧相关的上下文信息；$\pi$——任务提示。与需要梯度更新的监督方法不同，该范式更强调知识迁移、语义建模和推理组织能力，因此第二章的后续内容也将重点围绕“大模型如何理解视频语义”与“多智能体如何组织复杂决策”展开。

本文研究的重点是训练自由视频摘要，即在不重新训练专用视频摘要模型、不过度依赖大规模任务标注数据的前提下，直接利用大语言模型及多模态模型的通用知识、上下文推理能力与语义对齐能力完成摘要生成。这种设定与传统监督方法存在明显区别。监督式视频摘要方法通常需要以人工标注的重要性分数或摘要标签为目标进行参数学习，虽然在特定数据集上可能取得较高性能，但对标注成本、域内分布与训练资源具有较强依赖。相比之下，训练自由方法强调通过任务建模、提示工程、语义表示与推理组织来激发基础模型能力，从而提升方法的可迁移性、可解释性与工程部署灵活性。

进一步地，本文第三章与第四章虽然都属于训练自由视频摘要范式，但其建模视角存在差异。第三章侧重于将视频转换为多层级文本语义表示，再依据主题一致性与语义相关性生成重要性分数；第四章则将视频摘要视为一个需要规划、角色分工与多步决策的复杂推理任务，通过 Planner 与多个 Expert Agent 的协同实现多维度评分。二者共同体现出当前视频摘要研究从“视觉特征主导”向“语义推理主导”演进的趋势。

\xsection{面向层级语义视频摘要的大模型技术基础}{Large-model foundations for hierarchical semantic video summarization}

第三章提出的方法以大语言模型和多模态大语言模型为核心，通过“帧级描述生成-场景级语义聚合-视频级主题建模-语义相关性评分”的层级化语义链路，实现对原始视频内容的重要性估计。该方法能够成立，依赖于三个关键技术基础：其一，Transformer 为大语言模型提供统一的长程依赖建模能力；其二，提示工程与上下文学习使模型能够在不更新参数的条件下完成摘要、归纳与评分等任务迁移；其三，跨模态语义对齐使视觉内容能够被映射到统一的文本语义空间，从而支持帧、场景与视频主题之间的一致性比较\cite{vaswani2017attention,liu2023visual,ren2024timechat,song2024moviechat,he2024ma,yang2024qwen2}。

\xsubsection{Transformer 序列建模与表示学习机制}{Transformer-based sequence modeling and representation learning}

Transformer 是现代大语言模型的基础结构，其核心作用在于通过自注意力机制直接建模序列内部任意位置之间的依赖关系，从而突破传统循环神经网络在长序列处理中存在的梯度衰减与并行效率不足问题\cite{vaswani2017attention}。对于本文第三章而言，Transformer 的价值主要体现在两个方面：一方面，它能够在长文本上下文中同时考虑局部描述与全局主题之间的联系；另一方面，它为后续的场景级描述聚合、视频级摘要生成和重要性评分提供了统一的表示学习框架。

设输入词元序列的嵌入表示为 $X\in\mathbb{R}^{n\times d}$，其中 $n$ 为序列长度，$d$ 为特征维度。经线性映射后可得到查询矩阵 $Q$、键矩阵 $K$ 和值矩阵 $V$，则缩放点积注意力可表示为
\begin{equation}
\mathrm{Attention}(Q,K,V)=\mathrm{Softmax}\left(\frac{QK^{\top}}{\sqrt{d_k}}\right)V,
\end{equation}
式中，$Q$、$K$、$V$——分别为查询矩阵、键矩阵和值矩阵；$d_k$——键向量维度；$\mathrm{Attention}(\cdot)$——注意力输出。该公式说明，序列中每个位置的表示都会根据与其他位置的相关性进行加权更新，因此模型能够同时关注局部细节与整体语境。

为了从不同子空间捕获多粒度语义关系，Transformer 进一步采用多头注意力机制，其形式为
\begin{equation}
\mathrm{MultiHead}(Q,K,V)=\mathrm{Concat}(\mathrm{head}_1,\ldots,\mathrm{head}_h)W^O,
\end{equation}
式中，$\mathrm{head}_i=\mathrm{Attention}(QW_i^Q,KW_i^K,VW_i^V)$——第 $i$ 个注意力头输出；$h$——注意力头数；$W_i^Q$、$W_i^K$、$W_i^V$ 和 $W^O$——可学习参数矩阵。多头结构使模型能够分别关注实体关系、动作演变、上下文转折和主题线索等不同类型的语义模式，这对于视频文本描述中的多粒度信息建模尤为重要。

除自注意力之外，Transformer 还依赖位置编码与前馈网络来补足顺序信息和非线性表达能力。经典位置编码可写为
\begin{equation}
\mathrm{PE}(pos,2i)=\sin\left(\frac{pos}{10000^{2i/d}}\right), \qquad
\mathrm{PE}(pos,2i+1)=\cos\left(\frac{pos}{10000^{2i/d}}\right),
\end{equation}
式中，$pos$——词元位置索引；$i$——维度索引；$d$——嵌入维度。该编码使模型在并行计算条件下仍能感知序列先后顺序。随后，前馈网络进一步完成通道级非线性变换：
\begin{equation}
\mathrm{FFN}(x)=\sigma(xW_1+b_1)W_2+b_2,
\end{equation}
式中，$x$——输入表示；$W_1$、$W_2$——参数矩阵；$b_1$、$b_2$——偏置项；$\sigma(\cdot)$——非线性激活函数。由此，Transformer 可以同时建模顺序关系、全局依赖与非线性语义变换，为后续层级语义归纳提供完整的表示基础。

在第三章的方法中，场景级描述和视频级摘要本质上都可以视为对已有文本上下文的再组织与再压缩。Transformer 的自注意力结构使模型能够在同一上下文窗口内综合考虑不同帧描述之间的因果顺序、动作延续与语义重复，从而完成从细粒度视觉描述到高层语义概括的层级聚合。因此，Transformer 不仅是语言建模的底层结构，也是本文基于文本中介实现视频语义抽象的基础。

\xsubsection{基于提示学习的任务适配与推理机制}{Prompt-based task adaptation and reasoning}

提示工程是训练自由大语言模型应用中的关键技术，其基本思想是在不更新模型参数的前提下，通过构造任务指令、输入模板、角色约束和输出格式约束，引导模型在上下文中完成指定任务。与监督微调不同，上下文学习更加依赖提示设计质量，即模型是否能够从输入中准确识别任务目标、信息边界与响应形式。对于第三章的方法而言，场景摘要生成、视频主题概括以及场景重要性评分均不是单纯的开放式文本生成，而是带有明确目标约束的任务推理过程，因此提示工程直接影响生成结果的稳定性与可用性。

设输入查询为 $x$，上下文示例或背景信息为 $\mathcal{D}_c$，提示模板为 $\pi$，则模型输出可抽象表示为
\begin{equation}
y^{*}=\arg\max_{y} P_{\theta}(y\mid x,\mathcal{D}_c,\pi),
\end{equation}
式中，$y^{*}$——最优输出；$P_{\theta}(\cdot)$——参数固定的大语言模型定义的条件分布；$x$——当前查询；$\mathcal{D}_c$——上下文信息；$\pi$——提示模板。该表达式说明，在训练自由设定下，提示本身承担了“任务接口”的作用，模型性能不仅取决于预训练知识，也取决于输入组织方式。

从实现角度看，提示工程通常至少包含三类要素。第一类是任务描述，用于明确模型需要完成的是“总结”“比较”还是“评分”；第二类是上下文约束，用于给出视频级主题、场景级描述或历史已知信息，使模型能够在更完整的语境下进行判断；第三类是输出约束，用于限定模型返回分数、解释或结构化字段，以提升可解析性与可复用性。对于视频摘要任务，这种结构化提示尤其重要，因为模型不仅需要理解当前片段内容，还需要判断该片段相对于全局主题的代表性和独特性。

进一步地，提示学习可以按样本组织方式划分为零样本、少样本和链式推理等不同形式。零样本提示依赖模型已有知识直接完成任务，适合场景概括和开放式摘要；少样本提示通过少量示例显式说明输入输出对应关系，能够提高评分和格式化输出的稳定性；链式推理则通过中间分析步骤增强复杂判断的一致性。在训练自由视频摘要中，模型往往既需要理解视觉描述，又要输出可比较的连续分数，因此提示设计通常需要同时兼顾任务语义、推理路径和结果约束。

若将结构化输出映射为可用分数，可定义后处理函数
\begin{equation}
\hat{s}=\psi\big(\mathrm{LLM}(x,\mathcal{D}_c,\pi)\big), \qquad \psi(u)=\min\{1,\max(0,u)\},
\end{equation}
式中，$\hat{s}$——归一化后的重要性分数；$\mathrm{LLM}(\cdot)$——语言模型输出；$\psi(\cdot)$——分数裁剪函数。该形式说明，提示工程不仅影响文本内容本身，还影响输出能否顺利转化为后续计算所需的结构化变量，这一点在第三章的场景评分阶段和第四章的多 Agent 评分阶段都十分关键。

在第三章的方法中，提示工程具有两个直接作用。其一，它使大语言模型能够将原本开放式的视觉描述归纳任务，转化为具有明确主旨约束的层级摘要任务，即先从帧级描述中抽取场景语义，再从多个场景中归纳视频主旨。其二，它使场景评分过程由“单独判断内容是否重要”转变为“在全局主题约束下判断当前场景是否关键”，从而缓解局部显著但全局无关内容被误判为重要片段的问题。可以说，提示工程是第三章训练自由语义评分链路中的调度机制，也是大语言模型知识迁移到视频摘要任务中的关键接口。

\xsubsection{多模态语义表征与跨模态对齐机制}{Multimodal semantic representation and cross-modal alignment}

第三章方法的一个核心前提是：视频内容虽然最初以连续视觉信号形式存在，但在经过视觉语言模型描述后，可以被映射到统一的文本语义空间，并进一步与场景语义、视频主旨进行一致性比较。这一过程依赖跨模态语义对齐机制，即通过视觉编码、语言解码和共享语义空间映射，将视觉实体、动作关系和场景事件转化为可被语言模型处理的语义表示\cite{liu2023visual,zhang2023video,ren2024timechat,song2024moviechat,he2024ma,yang2024qwen2}。

设输入视频中采样得到的视觉特征序列为 $V=\{v_i\mid i=1,2,\ldots,N\}$，则跨模态投影过程可表示为
\begin{equation}
z_i=W_p v_i+b_p,
\end{equation}
式中，$v_i$——第 $i$ 个视觉标记的原始特征；$W_p$——投影矩阵；$b_p$——偏置项；$z_i$——映射后的语义嵌入。经过这一投影后，视觉内容被转换到与语言表征兼容的隐空间中，从而能够与文本提示或文本摘要共同参与推理。

在统一语义空间中，不同层级语义之间的相关性通常可以通过相似度函数进行度量。设帧级描述嵌入、场景级描述嵌入和视频级摘要嵌入分别为 $\mathbf{e}_f$、$\mathbf{e}_s$ 和 $\mathbf{e}_v$，则余弦相似度可定义为
\begin{equation}
\mathrm{sim}(\mathbf{a},\mathbf{b})=\frac{\mathbf{a}^{\top}\mathbf{b}}{\|\mathbf{a}\|_2\,\|\mathbf{b}\|_2},
\end{equation}
式中，$\mathbf{a}$ 和 $\mathbf{b}$——待比较的两个语义向量；$\|\cdot\|_2$——$L_2$ 范数；$\mathrm{sim}(\cdot,\cdot)$——余弦相似度。基于该定义，系统能够同时计算帧与所属场景之间的局部一致性，以及帧与全局主题之间的全局一致性，并据此实现分层重要性增强。

从层级语义建模的角度看，第三章的方法本质上构造了一条“视觉观测到语义判断”的逐级抽象链。若将帧级描述生成函数记为 $\Phi_f(\cdot)$，场景级语义聚合函数记为 $\Phi_s(\cdot)$，视频级主题归纳函数记为 $\Phi_v(\cdot)$，则有
\begin{equation}
\mathrm{Cap}_f(t)=\Phi_f(I_t), \qquad
\mathrm{Cap}_s(k)=\Phi_s\big(\{\mathrm{Cap}_f(t)\mid t\in S_k\}\big), \qquad
\mathrm{Cap}_v=\Phi_v\big(\{\mathrm{Cap}_s(k)\}_{k=1}^{N}\big),
\end{equation}
式中，$I_t$——第 $t$ 帧图像；$\mathrm{Cap}_f(t)$——帧级描述；$\mathrm{Cap}_s(k)$——第 $k$ 个场景的语义描述；$\mathrm{Cap}_v$——视频级全局摘要。该过程实现了由局部视觉内容到全局视频主题的逐步抽象，使后续重要性判断不再依赖单帧孤立信息，而是建立在层级语义一致性的基础之上。

在实际系统中，多模态对齐不仅承担“理解视频内容”的作用，还承担“把不同粒度信息组织到统一语义坐标系中”的作用。只有在这一前提下，系统才能将局部帧描述、场景摘要和视频主旨放到同一语义空间中进行比较，并据此实现主题约束下的重要性回传。因此，多模态语义表征与跨模态对齐既是第三章方法的输入接口，也是其评分机制成立的核心前提。

对于本文第三章而言，跨模态语义对齐的意义不只在于“生成描述”，更在于建立一种可比较、可传播的统一语义媒介。原始视频内容经过帧级描述后，被转化为语言可解释的局部语义单元；多个局部单元再经聚合得到场景级和视频级语义；最终，这些语义表示通过相似度度量反向作用于帧级评分，实现从高层主题约束到底层关键帧选择的闭环。因此，跨模态语义对齐是第三章层级化语义视频摘要方法得以成立的表示基础。

\xsection{面向协同推理视频摘要的智能体技术基础}{Agent foundations for collaborative reasoning in video summarization}

第四章提出的方法不再将视频摘要仅视为单一模型的直接打分问题，而是将其建模为一个需要规划、分工、交互与聚合的复杂推理过程。具体而言，该章采用 Planner 与多个 Expert Agent 组成的多智能体协同框架，对视频片段从叙事、视觉、情绪和信息等多个维度进行联合判断。该框架得以实现，主要依赖两项技术原理：一是基于大语言模型的 Agent 推理机制，使系统具备感知任务状态、生成决策与利用上下文记忆的能力；二是多 Agent 协作机制，使不同角色能够围绕同一目标形成互补分工与结果聚合。

\xsubsection{大语言模型驱动的智能体决策机制}{LLM-driven agent decision mechanism}

Agent 的核心思想是将大语言模型从“单轮响应器”扩展为“可持续决策体”。在这一框架下，模型不仅负责生成文本，还需要依据当前任务状态进行目标判断、行动选择、上下文更新和结果反馈。对于视频摘要任务而言，单次调用大语言模型虽然可以对一个片段给出粗略评价，但往往难以同时兼顾全局主题、局部上下文、历史片段与输出格式稳定性。因此，引入 Agent 机制后，系统可以把视频摘要过程组织为“观测-规划-评分-更新”的迭代推理链，以增强决策的一致性与可解释性。

设时刻 $t$ 的环境状态为 $s_t$，历史交互轨迹为 $\tau_t=\{(s_k,a_k,o_k)\mid k=1,2,\ldots,t-1\}$，记忆上下文为 $m_t$，则 Agent 在当前时刻的最优动作可表示为
\begin{equation}
a_t^{*}=\arg\max_{a} P_{\theta}(a\mid s_t,\tau_t,m_t),
\end{equation}
式中，$a_t^{*}$——时刻 $t$ 的最优动作；$P_{\theta}(\cdot)$——由大语言模型参数 $\theta$ 定义的策略分布；$s_t$——当前状态；$\tau_t$——历史交互轨迹；$m_t$——记忆上下文。该表达式说明，Agent 的决策不是静态地依赖当前输入，而是综合考虑了历史状态、先前动作及其反馈结果。

在执行动作后，系统还需基于反馈更新内部记忆，其状态更新过程可写为
\begin{equation}
m_{t+1}=f(m_t,s_t,a_t,o_t),
\end{equation}
式中，$m_{t+1}$——更新后的记忆状态；$f(\cdot)$——记忆更新函数；$a_t$——执行动作；$o_t$——环境反馈。该机制使 Agent 能够在处理后续片段时利用历史语义信息，避免把相邻内容完全视为相互独立的输入。对于第四章的方法，这种轻量级上下文记忆能够帮助模型识别叙事推进、重复片段和过渡片段，从而提升段级评分的时序一致性。

在本文第四章的具体框架中，Planner Agent 先基于全部片段描述生成视频主题、全局摘要与专家权重，随后再结合当前片段及其历史上下文输出 Planner 主分数。由此可见，Agent 推理机制在该章中承担了两类功能：一类是全局层面的规划，即确定评分策略与评价重点；另一类是局部层面的决策，即针对具体片段给出与当前上下文一致的分数判断。这使得视频摘要从“静态片段排序”转变为“带有状态的多步推理任务”。

若从系统执行角度进行归纳，Agent 决策通常可分解为感知、规划、执行和反馈四个阶段。感知阶段负责接收片段描述、历史记忆及任务约束；规划阶段负责决定后续关注重点和评价维度；执行阶段给出评分或调用外部工具；反馈阶段则根据中间结果更新内部状态。这样的闭环结构使 Agent 具备比单次调用更强的任务保持能力，也使视频摘要能够从“独立片段判断”转向“全局一致性驱动的动态判断”。

\xsubsection{多智能体协同建模与角色分工机制}{Multi-agent collaborative modeling and role specialization}

多 Agent 协作机制的基本思想是利用多个具有不同职责和能力偏好的智能体，共同完成单一智能体难以稳定处理的复杂任务。在视频摘要场景中，重要性并不只取决于视觉显著性，还可能受叙事推进、情绪表达和信息密度等因素影响。若仅依赖单个语言模型一次性做出综合判断，容易出现维度耦合不清、评分依据模糊以及对特定视频类型适配不足等问题。因此，第四章采用 Planner-Agent 与多 Expert-Agent 的协同模式，将复杂判断拆解为可解释的多角色决策过程。

设时刻 $t$ 的全局状态为 $s_t$，第 $i$ 个智能体的局部记忆为 $m_i$，协作上下文为 $c_t$，则该智能体的局部动作选择可表示为
\begin{equation}
a_i^{*}=\arg\max_{a_i} P_{\theta_i}(a_i\mid s_t,m_i,c_t),
\end{equation}
式中，$a_i^{*}$——第 $i$ 个智能体的最优局部动作；$P_{\theta_i}(\cdot)$——第 $i$ 个智能体的策略分布；$s_t$——全局状态；$m_i$——局部记忆；$c_t$——协作上下文。当各角色智能体分别完成局部推理后，其结果需要被进一步聚合为全局判断。设第 $i$ 个智能体输出的中间表示或评分为 $h_i$，则全局协同结果可写为
\begin{equation}
H_t=\sum_{i=1}^{N}\omega_i h_i, \qquad \omega_i=\frac{\exp(e_i)}{\sum_{j=1}^{N}\exp(e_j)},
\end{equation}
式中，$H_t$——时刻 $t$ 的全局协同结果；$\omega_i$——第 $i$ 个智能体的聚合权重；$e_i$——其重要性评分；$N$——参与协作的智能体数量。该形式体现了多 Agent 协作中“局部判断-加权聚合”的基本思想。

对于第四章所提出的方法，Planner Agent 负责从全局视角判断视频类型、总结主旨并分配各专家的重要性权重；Story Agent 聚焦叙事推进与事件转折；Visual Agent 聚焦画面变化与视觉显著性；Emotion Agent 聚焦情绪张力与情感表达；Information Agent 聚焦信息增量与知识密度。各 Expert Agent 的评分并非简单重复，而是分别从不同维度对当前片段的价值进行补充判断，最终由 Planner 给出的权重进行聚合。这样的设计一方面提升了评分依据的可解释性，另一方面也使系统能够根据视频类型动态调整多维评价侧重点。例如，教程类视频可提升信息维度权重，风景类视频则更依赖视觉维度判断。

若以第四章的评分过程为例，设 Planner 给出的主分数为 $p_k$，第 $m$ 个专家对片段 $S_k$ 的评分为 $e_k^{(m)}$，专家权重为 $w_m$，则片段最终重要性可写为
\begin{equation}
\hat{y}_k=p_k+\sum_{m\in\mathcal{M}} w_m e_k^{(m)}, \qquad \sum_{m\in\mathcal{M}} w_m = 1,
\end{equation}
式中，$\hat{y}_k$——第 $k$ 个片段的最终重要性分数；$\mathcal{M}$——专家集合；$w_m$——第 $m$ 个专家的权重。该公式清晰地体现出第四章方法的协同本质，即 Planner 负责建立全局评价框架，Experts 负责补充局部维度证据，二者共同构成最终决策。

从方法论角度看，多 Agent 协作机制在第四章中解决了三个关键问题。第一，它通过角色分工缓解了单一模型同时处理多评价维度时的推理负担；第二，它通过规划者与专家之间的层次化组织，提高了评分过程对全局目标的服从性；第三，它通过加权聚合与短期记忆注入，增强了不同片段间的时序一致性与跨片段关联建模能力。因此，多 Agent 协作并非简单地增加模型数量，而是为训练自由视频摘要提供了一种更具结构化、可解释性与可扩展性的推理组织方式。

\xsubsection{记忆增强与上下文状态维护机制}{Memory enhancement and contextual state maintenance}

对于长视频而言，单个片段的重要性往往依赖于其在前后叙事链中的位置。如果系统仅根据当前片段内容独立打分，就可能把重复出现的铺垫片段误判为多个关键时刻，或者忽略某些虽然视觉不显著、但在叙事上承前启后的过渡内容。因此，记忆增强机制是第四章多智能体推理框架中不可缺少的补充模块，其核心目标是在有限上下文预算下保留必要的历史语义线索，以增强当前决策的时序一致性。

设当前待评分片段为 $S_k$，历史窗口长度为 $h$，则其短期上下文记忆可表示为
\begin{equation}
\mathcal{M}_k^{(h)} = \{c_{k-h},c_{k-h+1},\ldots,c_{k-1}\},
\end{equation}
式中，$\mathcal{M}_k^{(h)}$——片段 $S_k$ 的历史上下文记忆；$c_i$——第 $i$ 个片段的文本描述；$h$——历史窗口长度。通过将这些历史描述与当前片段共同输入 Planner 或 Expert Agent，系统能够判断当前内容究竟是在推进主线、重复已有信息，还是构成语义过渡。

记忆机制的价值不仅在于提供更多上下文，还在于改变评分依据的语义粒度。没有记忆时，模型更容易基于片段本身的局部显著性做出判断；引入记忆后，模型则能够从“相对于前文是否新增信息”“是否形成情绪转折”“是否重复已出现事件”等角度进行更高层次的比较。对于第四章的方法而言，这一机制直接支撑了 Story Agent、Information Agent 等角色在连续片段之间建立差异化判断，也使 Planner 能够根据局部历史动态调整评价尺度。

\xsection{视频摘要评价指标与评测协议}{Evaluation metrics and protocols for video summarization}

由于本文第三章与第四章的实验均在相同的公开数据集上展开，且采用统一的评估体系，本节对视频摘要任务中使用的核心评估指标与评测协议进行集中介绍，以避免在后续各章中重复阐述。需要指出的是，视频摘要评价既涉及最终选出的二值摘要是否与人工摘要一致，也涉及模型输出的连续重要性分数能否保持与人工标注相近的相对排序。因此，本文同时采用集合重合度指标与排序一致性指标对模型进行综合评估。

本文采用视频摘要领域最为通用的 $F_1$-Score 作为核心定量评估指标，用以衡量模型生成摘要与人工标注摘要之间的一致性。

视频摘要任务的最终输出为一个与原视频等长的二值向量，其中值为 1 的位置表示对应帧被选入摘要，值为 0 则表示被排除。设模型生成的二值摘要向量为 $\mathbf{S}_{\mathrm{sys}}$，某一标注者给出的参考摘要向量为 $\mathbf{S}_{\mathrm{gt}}$，二者的元素取值均属于 $\{0,1\}$。在此基础上，精确率（Precision）与召回率（Recall）分别定义为
\begin{equation}
  P = \frac{|\mathbf{S}_{\mathrm{sys}}\cap \mathbf{S}_{\mathrm{gt}}|}{|\mathbf{S}_{\mathrm{sys}}|}, \qquad
  R = \frac{|\mathbf{S}_{\mathrm{sys}}\cap \mathbf{S}_{\mathrm{gt}}|}{|\mathbf{S}_{\mathrm{gt}}|},
\end{equation}
式中，$P$——精确率；$R$——召回率；$\mathbf{S}_{\mathrm{sys}}$——模型生成的二值摘要向量；$\mathbf{S}_{\mathrm{gt}}$——人工参考摘要向量；$|\cdot|$——集合的基数。
其中 $|\cdot|$ 表示集合的基数，$\mathbf{S}_{\mathrm{sys}}\cap \mathbf{S}_{\mathrm{gt}}$ 为模型所选帧与标注所选帧的交集。精确率 $P$ 反映模型所选帧中实际属于参考摘要的比例，召回率 $R$ 则反映参考摘要中被模型成功召回的比例。二者的调和平均即为 $F_1$ 分数：
\begin{equation}
  F_{1}= \frac{2PR}{P + R}.
\end{equation}
式中，$F_1$——精确率与召回率的调和平均指标；$P$——精确率；$R$——召回率。

由于两个数据集均包含多位标注者的独立标注，单一参考摘要无法完整反映视频摘要任务固有的主观性。针对不同数据集的标注特点，领域内分别采用两种评估协议。对于 SumMe 数据集，将模型摘要与每位标注者的参考摘要逐一计算 $F_1$ 分数后取最大值，该协议称为 Max $F_1$-Score：
\begin{equation}
  F_1^{\max}(v) = \max_{u \in \mathcal{U}_v} F_1(\mathbf{S}_{\mathrm{sys}}^{(v)},\, \mathbf{S}_{\mathrm{gt}}^{(v,u)}),
\end{equation}
式中，$F_1^{\max}(v)$——视频 $v$ 在 SumMe 协议下的最大 $F_1$ 分数；$\mathcal{U}_v$——视频 $v$ 的标注者集合；$\mathbf{S}_{\mathrm{sys}}^{(v)}$——视频 $v$ 的模型摘要；$\mathbf{S}_{\mathrm{gt}}^{(v,u)}$——标注者 $u$ 给出的参考摘要。
其中 $\mathcal{U}_v$ 为视频 $v$ 的标注者集合，$\mathbf{S}_{\mathrm{gt}}^{(v,u)}$ 为标注者 $u$ 给出的参考摘要。对于 TVSum 数据集，先将多位标注者的逐帧重要性评分取均值并二值化为参考摘要，再与模型输出计算 $F_1$ 分数，该协议称为 Avg $F_1$-Score：
\begin{equation}
  F_1^{\mathrm{avg}}(v) = F_1\!\big(\mathbf{S}_{\mathrm{sys}}^{(v)},\, \bar{\mathbf{S}}_{\mathrm{gt}}^{(v)}\big), \qquad \bar{\mathbf{S}}_{\mathrm{gt}}^{(v)} = \mathrm{Binarize}\!\Big(\frac{1}{|\mathcal{U}_v|}\sum_{u \in \mathcal{U}_v}\mathbf{r}^{(v,u)}\Big)
\end{equation}
式中，$F_1^{\mathrm{avg}}(v)$——视频 $v$ 在 TVSum 协议下的平均 $F_1$ 分数；$\bar{\mathbf{S}}_{\mathrm{gt}}^{(v)}$——由多位标注者评分聚合得到的参考摘要；$\mathbf{r}^{(v,u)}$——标注者 $u$ 给出的帧级重要性评分向量；$\mathrm{Binarize}(\cdot)$——二值化操作。
  其中 $\mathbf{r}^{(v,u)}$ 为标注者 $u$ 给出的帧级重要性评分向量，$\mathrm{Binarize}(\cdot)$ 为按阈值二值化操作。最终，在各数据集上对所有视频的对应 $F_1$ 指标取均值即得总体评估分数。本文严格遵循上述协议——SumMe 报告 Max $F_1$-Score，TVSum 报告 Avg $F_1$-Score——以确保与现有方法具有公平可比性。

  除二值摘要层面的 $F_1$-Score 外，本文还关注模型输出的连续重要性分数与人工评分序列之间的排序一致性，并采用 Spearman 秩相关系数 $\rho$ 与 Kendall 秩相关系数 $\tau$ 作为辅助评估指标。设模型预测分数序列为 $\mathbf{s}=(s_1,\ldots,s_T)$，人工评分序列为 $\mathbf{r}=(r_1,\ldots,r_T)$，则 Spearman 秩相关系数可写为
  \begin{equation}
  \rho = 1 - \frac{6\sum_{t=1}^{T} d_t^2}{T(T^2-1)},
  \end{equation}
  式中，$d_t$——第 $t$ 个位置在两个排序中的秩差；$T$——序列长度；$\rho$——取值范围为 $[-1,1]$ 的 Spearman 秩相关系数。当 $\rho$ 越接近 1 时，说明模型对相对重要性顺序的保持能力越强。相应地，Kendall 秩相关系数定义为
  \begin{equation}
  \tau = \frac{N_c - N_d}{\frac{1}{2}T(T-1)},
  \end{equation}
  式中，$N_c$——一致样本对数量；$N_d$——不一致样本对数量；$\tau$——取值于 $[-1,1]$ 的 Kendall 秩相关系数。与 $F_1$-Score 强调最终摘要重合度不同，$\rho$ 和 $\tau$ 更关注连续分数层面的排序合理性，因此能够补充刻画模型在“先排序、后选取”流程中的质量。

  在实际评测过程中，模型通常首先输出帧级或片段级连续分数，然后在统一预算约束下生成最终摘要。设摘要长度比例为 $\sigma$，原视频总帧数为 $T$，则预算约束可写为 $B=\sigma T$。在公开视频摘要研究中，常见设置是令摘要长度不超过原视频时长的 15\%。对片段级方法而言，还需先将片段分数映射回对应帧区间，再通过平滑、二值化和背包优化等步骤生成最终摘要序列。只有在统一预算和统一协议下进行评测，不同方法的实验结果才具有可比性。正因如此，本文在第三章和第四章中均严格遵循相同的评测流程，以保证实验结论的客观性和一致性。

